% Generated from validated Special post-draft synthesis. Do not hand-edit.
\section*{Monthly Signals}
\addcontentsline{toc}{section}{Monthly Signals}
\smallskip
\noindent \textbf{Frontier Models} 性能順位だけでなく、open-weight、multimodal、API lifecycle、computer-use、coding/agent統合という提供形態の違いが競争軸として前景化した。 \hfill{\footnotesize p.~\pageref{pkg:frontier-models-june}}
\par
\smallskip
\noindent \textbf{長時間Agentの実行と安全} 長時間taskではmodel能力に加えて、scheduling、persistent memory、context transformation、per-call authorization、複数toolをまたぐtrust boundaryがsystem設計の中心課題になった。 \hfill{\footnotesize p.~\pageref{pkg:agents-memory-june} / p.~\pageref{pkg:agent-safety-security-june}}
\par
\smallskip
\noindent \textbf{Inference \& Serving} custom accelerator、kernel、parallelism、serving runtimeが同時に更新され、推論最適化が単一componentではなくstack全体の問題として見えるようになった。 \hfill{\footnotesize p.~\pageref{pkg:inference-serving-june}}
\par
\smallskip
\noindent \textbf{AI for Science} near-autonomous chemistry workflowは、AIが提案だけでなく実験工程を含む長時間workflowへ入り始めたことと、human-in-the-loopの責任境界を同時に示した。 \hfill{\footnotesize p.~\pageref{pkg:ai-for-science-june}}
\par
\smallskip
\noindent \textbf{Paper Watch} execution-time safety kernel、context rot、memory evaluation confoundが、長時間Agentを制御し正しく測るための研究課題として立ち上がった。 \hfill{\footnotesize p.~\pageref{pkg:paper-watch-june}}
\par
\medskip
\begin{claimboundary}[Retrospective scope]
本号は2026年7月を、後日確認可能になった一次情報も用いて再構成するRetrospective Specialである。Coverage Periodと制作時点を同一視せず、vendor / project / author claimの境界は各記事で明示する。
\end{claimboundary}
\medskip
\tableofcontents
